\begin{lemma}[Sobolev chain rule for the smooth regularized absolute value]
For $\delta > 0$ let $q_\delta(t) = \sqrt{t^2+\delta^2} - \delta$.  Let
$u \in W^{1,2}(\gamma_n)$ in the sense of \noderef{SobolevW12} and let
$w \in L^2(\gamma_n)$ satisfy $w = q_\delta \circ u$ $\gamma_n$-almost
everywhere.  Then $w \in W^{1,2}(\gamma_n)$ and
\[
  Dw = \frac{u}{\sqrt{u^2+\delta^2}}\; Du
  \qquad \gamma_n\text{-a.e.},
\]
with $D$ the energy-normalized weak gradient of \noderef{GaussD}.  No regularity
of the set $\{u>0\}$ or of its boundary is used; the statement concerns the local
weak gradient only, not any locality of the nonlocal operator $B$.
\end{lemma}
\begin{proof}
\emph{Step 0: properties of $q_\delta$ and notation.}
$q_\delta$ is smooth on $\mathbb{R}$ with
$q_\delta'(t) = t/\sqrt{t^2+\delta^2}$ and
$q_\delta''(t) = \delta^2 (t^2+\delta^2)^{-3/2}$.  Hence
$|q_\delta'| \leq 1$ everywhere, $q_\delta$ is $1$-Lipschitz, $q_\delta(0) = 0$,
and $q_\delta'$ is Lipschitz with constant
$\sup|q_\delta''| = \delta^2\cdot\delta^{-3} = 1/\delta$.  Set
\[
  G = \sqrt2\, Du \in L^2(\gamma_n;\mathbb{R}^n),
  \qquad
  \Theta(x) = q_\delta'\bigl(u(x)\bigr) = \frac{u(x)}{\sqrt{u(x)^2+\delta^2}} ,
\]
so $G$ is a weak Euclidean gradient of $u$ by \noderef{SobolevW12} and
\noderef{GaussD}, and $|\Theta| \leq 1$ pointwise.  Then
$\Theta\, G \in L^2(\gamma_n;\mathbb{R}^n)$, since $|\Theta G| \leq |G|$
pointwise.  Our goal is to show that $\Theta\,G$ is a weak Euclidean gradient of
$w$; the conclusion then follows, because $\Theta G = \sqrt2\,(\Theta\, Du)$ with
$\Theta\,Du \in L^2(\gamma_n;\mathbb{R}^n)$ gives $w \in W^{1,2}(\gamma_n)$ by
\noderef{SobolevW12}, and \noderef{GaussD} together with
\noderef{GaussWeakGradUnique} identifies $Dw = \Theta\,Du$, which is the asserted
formula.

Since $w = q_\delta\circ u$ $\gamma_n$-a.e., and $\gamma_n$ and Lebesgue measure
have the same null sets (the density of \noderef{GaussMeasure} is continuous and
strictly positive), we may and do replace $w$ by $q_\delta\circ u$ throughout:
changing a function on a null set changes neither side of the defining identity
of \noderef{HasGaussWeakGrad}.

\emph{Step 1: reduction to a fixed test function and local mollification.}
Fix $j \in \{1,\ldots,n\}$ and $\varphi \in C_c^\infty(\mathbb{R}^n)$; we must show
\begin{equation*}
  \int_{\mathbb{R}^n} q_\delta(u)\, \partial_j \varphi\, dx
    = -\int_{\mathbb{R}^n} \Theta\, G_j\, \varphi \, dx . \tag{1}
\end{equation*}
Let $K = \operatorname{supp}\varphi$, a compact subset of $\mathbb{R}^n$, let
$\lambda$ denote Lebesgue measure, and put
$K_1 = \{x : \operatorname{dist}(x,K) \leq 1\}$, which is compact.

By \noderef{GaussL2LocallyIntegrable}, applied in turn to $u$ and to each
component $G_j$ of $G$ -- all of these lie in $L^2(\gamma_n)$, $u$ by hypothesis
and $G_j$ because $G \in L^2(\gamma_n;\mathbb{R}^n)$ -- each of
$u, G_1, \ldots, G_n$ is locally $\lambda$-integrable and lies in
$L^2(K',\lambda)$ for every compact $K' \subseteq \mathbb{R}^n$.

Fix a mollifier family $(\rho_\varepsilon)_{\varepsilon>0}$: each
$\rho_\varepsilon$ smooth, nonnegative, supported in $\{|t| \leq \varepsilon\}$,
with $\int\rho_\varepsilon\,d\lambda = 1$; for instance
$\rho_\varepsilon(t) = \varepsilon^{-n}\rho(t/\varepsilon)$ for a fixed smooth
nonnegative $\rho$ supported in the unit ball with $\int\rho\,d\lambda = 1$.  Put
$u_\varepsilon = \rho_\varepsilon * u$, that is
$u_\varepsilon(x) = \int \rho_\varepsilon(t)\,u(x-t)\,d\lambda(t)$.  Since $u$ is
locally $\lambda$-integrable and $\rho_\varepsilon \in C_c^\infty(\mathbb{R}^n)$,
the mollification item of shared context gives that $u_\varepsilon$ is defined at
every point and is $C^\infty$ on $\mathbb{R}^n$.

\emph{The gradient identity.}  For every $x \in \mathbb{R}^n$, every $j$ and every
$\varepsilon > 0$,
\begin{equation*}
  \partial_j u_\varepsilon(x) = (\rho_\varepsilon * G_j)(x) . \tag{1a}
\end{equation*}
Indeed, fix $x$ and put $\psi(y) = \rho_\varepsilon(x-y)$, so that
$\psi \in C_c^\infty(\mathbb{R}^n)$ and
$\partial_{y_j}\psi(y) = -(\partial_j\rho_\varepsilon)(x-y)$.  Differentiating the
convolution in $x$ moves the derivative onto the smooth factor,
$\partial_j u_\varepsilon = (\partial_j\rho_\varepsilon) * u$, and rewriting that
integral by the change of variable $y = x - t$, which preserves $\lambda$ by
translation and reflection invariance of Lebesgue measure (shared context), gives
\[
  \partial_j u_\varepsilon(x)
    = \int u(y)\,(\partial_j\rho_\varepsilon)(x-y)\,d\lambda(y)
    = -\int u\,\partial_j\psi\, d\lambda
    = \int G_j\,\psi\, d\lambda
    = (\rho_\varepsilon * G_j)(x) ,
\]
the third equality being the defining identity of \noderef{HasGaussWeakGrad} for
the pair $(u,G)$ with test function $\psi$, and the fourth the same change of
variable read backwards.

\emph{The convergences.}  Apply \noderef{MollifierL2Local} twice with the compact
set $K_1$ and the fixed mollifier family: once to $f = u$ and once to $f = G_j$.
Its hypothesis, that $f$ lie in $L^2(K',\lambda)$ for every compact $K'$, is
exactly what \noderef{GaussL2LocallyIntegrable} supplied above for both
functions.  Hence, as $\varepsilon \downarrow 0$ through positive values,
\[
  u_\varepsilon \to u \quad\text{in } L^2(K_1,\lambda),
  \qquad
  \partial_j u_\varepsilon \to G_j \quad\text{in } L^2(K_1,\lambda),
\]
the second statement using $(1a)$ to identify $\partial_j u_\varepsilon$ with
$\rho_\varepsilon * G_j$.  Since $K \subseteq K_1$, the same convergences hold in
$L^2(K,\lambda)$, because $\|w\|_{L^2(K,\lambda)} \leq \|w\|_{L^2(K_1,\lambda)}$
for every $w$.

\emph{Step 2: the classical chain rule for the mollifications.}
Each $u_\varepsilon$ is smooth on a neighbourhood of $K$ and $q_\delta$ is smooth
on $\mathbb{R}$, so the classical chain rule gives, pointwise on $K$,
\[
  \partial_j\bigl(q_\delta(u_\varepsilon)\bigr) = q_\delta'(u_\varepsilon)\,\partial_j u_\varepsilon .
\]
Both sides are continuous on a neighbourhood of $K$, and $\varphi$ is smooth with
support in $K$, so ordinary integration by parts gives
\begin{equation*}
  \int q_\delta(u_\varepsilon)\,\partial_j\varphi \, dx
    = -\int q_\delta'(u_\varepsilon)\,\partial_j u_\varepsilon\; \varphi\, dx . \tag{2}
\end{equation*}

\emph{Step 3: passing to the limit in (2).}
Left-hand side: $|q_\delta(u_\varepsilon) - q_\delta(u)| \leq |u_\varepsilon - u|$
pointwise, because $q_\delta$ is $1$-Lipschitz; hence
$q_\delta(u_\varepsilon) \to q_\delta(u)$ in $L^1(K)$ by Step 1, and since
$\partial_j\varphi$ is bounded and supported in $K$,
$\int q_\delta(u_\varepsilon)\partial_j\varphi\,dx \to
\int q_\delta(u)\,\partial_j\varphi\, dx$.

Right-hand side: write
\[
  q_\delta'(u_\varepsilon)\,\partial_j u_\varepsilon - \Theta\, G_j
    = q_\delta'(u_\varepsilon)\bigl(\partial_j u_\varepsilon - G_j\bigr)
      + \bigl(q_\delta'(u_\varepsilon) - q_\delta'(u)\bigr) G_j .
\]
The first term tends to $0$ in $L^1(K)$ because $|q_\delta'| \leq 1$ and
$\partial_j u_\varepsilon \to G_j$ in $L^2(K,\lambda)$ and hence in
$L^1(K,\lambda)$ (Step 1, with Cauchy--Schwarz on the finite-measure set $K$).  For the second term,
$q_\delta'$ is $1/\delta$-Lipschitz by Step 0, so
$|q_\delta'(u_\varepsilon) - q_\delta'(u)| \leq |u_\varepsilon - u|/\delta$, and
by Cauchy--Schwarz on $K$,
\[
  \int_K \bigl|q_\delta'(u_\varepsilon) - q_\delta'(u)\bigr|\,|G_j|\, dx
    \leq \frac{1}{\delta}\,\|u_\varepsilon - u\|_{L^2(K)}\,\|G_j\|_{L^2(K)}
    \longrightarrow 0 .
\]
Hence $\int q_\delta'(u_\varepsilon)\partial_ju_\varepsilon\,\varphi\,dx \to
\int \Theta\, G_j\,\varphi\, dx$, $\varphi$ being bounded with support in $K$.

Letting $\varepsilon \downarrow 0$ in (2) therefore yields (1).

\emph{Step 4: conclusion.}
Since $j$ and $\varphi$ were arbitrary, (1) says exactly that $\Theta\,G$ is a
weak Euclidean gradient of $q_\delta\circ u = w$ in the sense of
\noderef{HasGaussWeakGrad}.  By Step 0 this gives $w \in W^{1,2}(\gamma_n)$ and
$Dw = \Theta\, Du = \dfrac{u}{\sqrt{u^2+\delta^2}}\,Du$ $\gamma_n$-a.e.
\end{proof}
